\chapter{Introduction} \label{chap:intro}
\section{Contextualization}
%MIR intro
Music Information Retrieval (MIR) is a research field that develops computational methods to analyze, interpret, and retrieve musical information from performance recordings and symbolic representations. MIR bridges a wide range of disciplines, including musicology, signal processing, machine learning, and human-computer interaction. Because MIR is a data-intensive field, visualization plays a crucial role in its research. The way information is selected and represented shapes the conclusions we draw and guides further progress.

%The MIR visualization challenge
The challenge of visualization in MIR is particularly significant because music is an inherently abstract and subjective medium. Translating music into the visual domain is a problem that spans millennia; throughout history, multiple systems of musical notation have been developed in attempts to capture sound in written form, yet this remains an evolving challenge. For MIR purposes, this means that musical data comes in a variety of forms, including audio recordings, symbolic scores, and annotations of expressive qualities. Each of these requires a different visual approach. An audio visualization such as a waveform or spectrogram reveals how a performance sounds, while a musical score conveys how a piece should be performed; both serve distinct and important purposes for different research tasks.

\pagebreak

\section{Motivation}
% MIR Visualization Problem Motivation intro 
MIR researchers often need to inspect algorithmic outputs alongside audio representations and symbolic scores in a temporally coherent manner. This kind of multimodal, time-aligned visualization is central to both the evaluation of MIR algorithms and the development of annotation workflows. Audio-to-score alignment, the task of synchronizing a performance recording with its corresponding symbolic score, is a key enabler of this type of visualization, as it provides the temporal anchor needed to align heterogeneous data sources into a unified view.

% Focus on BT/BA
This thesis was motivated by two closely related challenges in beat-tracking (BT) and beat-annotation (BA) research. Contemporary BT systems are predominantly based on deep learning. Consequently, the field relies heavily on tools that support annotation and evaluation tasks, both of which are essential for advancing BT research.

Annotation tasks, particularly beat annotation, remain laborious and time-consuming. Existing graphical user interface (GUI) tools such as Sonic Visualiser leave a clear gap in this area: integrating recent state-of-the-art BT algorithms is difficult, and BA-specific features have seen little meaningful innovation. This disconnect between annotation workflows and the algorithms they are intended to support is a concrete obstacle to progress in both domains.

% Evaluation Challenge
GUI tools present a similar limitation for evaluation: they often lack methods for inspecting algorithmic behaviour in depth. Better tools for evaluating and comparing models with one another are essential for improving them.

% Stating the SV problem
Although solutions like Sonic Visualiser \cite{cannam2010SonicVisualiserOpen} are powerful and widely used, their development cannot keep pace with the needs of modern MIR research. Integrating new algorithms or producing specialized visualizations is often impractical within them. A particularly underserved case is the visualization of score-aligned data: combining audio and data representations with a musical score typically involves manual image manipulation, and no unified methodology is widely used in the community.

For these reasons, researchers frequently resort to \textit{ad hoc} solutions that lack reproducibility and are difficult to share or extend. This fragmentation creates a significant obstacle to the collaboration on which MIR research depends. By contributing a flexible, extensible, and community-oriented visualization tool built on established libraries, this thesis aims not only to address an immediate practical gap but also to establish a foundation for a framework that can support the growing complexity and interdisciplinarity of MIR research.

%Bridge to RQ and objectives
All of the concerns described in this section provide a basis for the objectives and research questions presented in the following section.
\pagebreak
\section{Objectives, Research Questions and Contributions}
This thesis is guided by three objectives, which in turn lead to three research questions and a set of concrete contributions.

\begin{itemize}
    \item \textbf{Objective 1}: Develop a reusable Python framework for composing score-aligned MIR visualizations based on an OO layer-composition model.
    \item \textbf{Objective 2}: Assess whether the framework's modular layers and scriptable composition support the generation, modification, and regeneration of score-aligned visualizations without changes to the composition mechanism.
    \item \textbf{Objective 3}: Assess how score-aligned visualizations can provide qualitative visual context for beat annotation and BT evaluation workflows.
\end{itemize}

\begin{itemize}
    \item \textbf{RQ1}: How can heterogeneous MIR visualizations be combined into a common score-aligned representation?
    \item \textbf{RQ2}: To what extent do LayerIt's modular layers and scriptable composition support generating, modifying, and regenerating score-aligned visualizations, including the demonstrated addition of a new time-based layer without modifying the composition mechanism?
    \item \textbf{RQ3}: How can score-aligned visualizations support qualitative inspection of beat-tracking outputs and annotation data on a common time axis in the demonstrated case studies, subject to the limitations of the available score-performance alignment?
\end{itemize}

The main contributions of this thesis are as follows:
\begin{itemize}
    \item The LayerIt framework for composing score-aligned MIR visualizations through an OO layer-composition model.
    \item The release of LayerIt as an open-source Python framework.
    \item The open-source repository containing examples of the visualizations developed as part of this work.
    \item A Late-Breaking Demo paper submitted to ISMIR 2026.
\end{itemize}

The primary contribution is the layer-composition framework and the figures it enables, rather than the alignment method itself, which is provided by ScoreWarp \cite{goeblLetsScoreWarpAgain2025a}.

\section{Developed Project} 
To address these objectives and research questions, we developed LayerIt, a Python-based object-oriented (OO) framework for composing score-aligned MIR visualizations. LayerIt integrates established MIR libraries and tools into a Scalable Vector Graphics (SVG)-based workflow in which visual elements are treated as independent layers that can be composed, reordered, and combined into custom figures.

At its core, LayerIt supports the integration of symbolic scores with audio and algorithmic data in a single view. By relying on ScoreWarp for temporal alignment, it makes it possible to overlay spectrograms, onset information, and beat-tracking outputs on score representations within the same visualization. The result is a flexible layer-composition model that supports the visualization workflows developed in this thesis and provides a foundation for future MIR visualization workflows.

\section{Thesis Structure}

The remainder of this thesis is organized as follows:
\begin{itemize}
    \item \textbf{Chapter~\ref{chap:sota}} provides a literature review covering the main audio visualization types, existing GUI and code-based visualization tools, and audio-to-score alignment.
    \item \textbf{Chapter~\ref{chap:Methodology}} presents the adopted methodology and describes the five iterative phases that culminated in the development of LayerIt.
    \item \textbf{Chapter~\ref{chap:DesignAndImplementation}} describes the problem addressed by LayerIt, its design and implementation, and the ``layers within layers'' philosophy that underpins its architecture and use.
    \item \textbf{Chapter~\ref{chap:Validation}} evaluates the framework against the research questions through the generated figures, a qualitative workflow comparison, and a discussion of its limitations.
    \item  \textbf{Chapter~\ref{chap:conclusion}} reflects on the achievements of the thesis and outlines directions for future work.
\end{itemize} 